\begin{trapbox}
	\begin{description}[style=nextline, leftmargin=2.8em, labelsep=0.5em, itemsep=0.35em]
		\item[\textbf{\textcolor{CTrap}{易错点一（条件混淆）}}] 误将对称轴条件当作 $\omega x+\varphi = k\pi$，对称中心条件当作 $\omega x+\varphi = k\pi+\frac{\pi}{2}$，导致结论完全颠倒。
		\item[\textbf{\textcolor{CTrap}{正确理解（正确条件）：}}] 对称轴要求相位等于 $\frac{\pi}{2}+k\pi$，对称中心要求相位等于 $k\pi$。
		\item[\textbf{\textcolor{CTrap}{错因分析（记忆错误）：}}] 对正弦函数取最值和零点的条件记忆不清，混淆了对称轴和对称中心的特征。
	\end{description}

	\medskip
	\begin{description}[style=nextline, leftmargin=2.8em, labelsep=0.5em, itemsep=0.35em]
		\item[\textbf{\textcolor{CTrap}{易错点二（漏除系数）}}] 在解 $x$ 时，直接写出 $x = k\pi+\frac{\pi}{2}-\varphi$ 或 $x = k\pi-\varphi$，忘记除以 $\omega$。
		\item[\textbf{\textcolor{CTrap}{正确理解（系数处理）：}}] 必须将整个表达式除以 $\omega$ 才能解出 $x$，注意 $\omega\neq 0$。
		\item[\textbf{\textcolor{CTrap}{错因分析（解法错误）：}}] 对整体代换后解方程的操作不熟练，忽略系数 $\omega$ 的作用。
	\end{description}

	\medskip
	\begin{description}[style=nextline, leftmargin=2.8em, labelsep=0.5em, itemsep=0.35em]
		\item[\textbf{\textcolor{CTrap}{易错点三（忽略整数）}}] 只写出 $k=0$ 对应的一个对称轴或对称中心，误认为只有一个。
		\item[\textbf{\textcolor{CTrap}{正确理解（整数标记）：}}] $k$ 必须取遍所有整数 $\mathbb{Z}$，表示无穷多条对称轴和无穷多个对称中心。
		\item[\textbf{\textcolor{CTrap}{错因分析（周期忽略）：}}] 对正弦函数的周期性理解不足，未意识到对称轴和对称中心也是周期性出现的。
	\end{description}
\end{trapbox}
